
\documentclass{article} % For LaTeX2e
\usepackage{iclr2026_conference,times}

% Optional math commands from https://github.com/goodfeli/dlbook_notation.
\input{math_commands.tex}

\usepackage{hyperref}
\usepackage{url}
\usepackage{booktabs}
\usepackage{multirow}


\title{Dynamic Test-Time Memory for Long-Horizon Alignment in Digital Psychotherapy}

% Authors must not appear in the submitted version. They should be hidden
% as long as the \iclrfinalcopy macro remains commented out below.
% Non-anonymous submissions will be rejected without review.

\author{Antiquus S.~Hippocampus, Natalia Cerebro \& Amelie P. Amygdale \thanks{ Use footnote for providing further information
about author (webpage, alternative address)---\emph{not} for acknowledging
funding agencies.  Funding acknowledgements go at the end of the paper.} \\
Department of Computer Science\\
Cranberry-Lemon University\\
Pittsburgh, PA 15213, USA \\
\texttt{\{hippo,brain,jen\}@cs.cranberry-lemon.edu} \\
\And
Ji Q. Ren \& Yevgeny LeNet \\
Department of Computational Neuroscience \\
University of the Witwatersrand \\
Joburg, South Africa \\
\texttt{\{robot,net\}@wits.ac.za} \\
\AND
Coauthor \\
Affiliation \\
Address \\
\texttt{email}
}

% The \author macro works with any number of authors. There are two commands
% used to separate the names and addresses of multiple authors: \And and \AND.
%
% Using \And between authors leaves it to \LaTeX{} to determine where to break
% the lines. Using \AND forces a linebreak at that point. So, if \LaTeX{}
% puts 3 of 4 authors names on the first line, and the last on the second
% line, try using \AND instead of \And before the third author name.

\newcommand{\fix}{\marginpar{FIX}}
\newcommand{\new}{\marginpar{NEW}}

%\iclrfinalcopy % Uncomment for camera-ready version, but NOT for submission.
\begin{document}


\maketitle

\begin{abstract}
Large Language Models (LLMs) often struggle with long-horizon coherence and clinical alignment in multi-session therapy. We propose a Dynamic Test-Time Memory (DTM) framework that utilizes a Generator-Extractor loop to facilitate inference-time continual learning. By evaluating on synthetic therapy transcripts using CBT-specific metrics and Gemini-as-a-judge, we demonstrate that DTM significantly improves persona consistency and clinical adherence compared to existing memory-augmented baselines like Mem0.
\end{abstract}

\section{Introduction}

Maintaining the therapeutic thread across multiple sessions is a major challenge for AI agents. We introduce the problem of "clinical drift" and propose our Dynamic Test-Time Memory approach as a solution for alignment in high-stakes mental health domains.

\begin{itemize}
    \item The Problem: LLMs drift in long-term counseling. Standard context windows lose the ``clinical thread," and RAG (like Mem0) often retrieves fragmented facts rather than therapeutic trajectories.
    \item The Goal: Maintaining ``Clinical Alignment" (adherence to CBT) and ``Persona Consistency" across multiple sessions/timelines.
    \item The Contribution: 1. Introduction of a Dynamic Test-Time Memory (DTM) architecture.
    2. A novel evaluation framework using CBT Scoring and Gemini-as-a-Judge on synthetic human-patient transcripts.
    3. Empirical evidence that a Generator-Extractor loop facilitates a form of ``Inference-time Continual Learning."

\end{itemize}

\section{Related Work} - Joel and Alexis
\label{gen_inst}

\subsection{Memory-Augmented LLMs}
Though capable of storing vast amounts of parametric knowledge after pre-training, LLMs still struggle with hallucinating plausible but unfactual information due to their limited context windows and degrading performance as context size increases. To address this, a number of works have been proposed in recent years. Retrieval Augmented Generation (Piktus et al. 2021) first introduced the idea of cosine similarity between query and memory chunk vector embeddings, eliminating the need to pass in entire contexts into LLMs to answer specific questions.  LongMem (Wang et al. 2023) introduced a novel decoupled network architecture with the original LLM frozen as a memory encoder, and an adaptive residual side-network as a memory retriever and reader. (Lou et al. 2025) combined a BERT-based classifier with an LSTM memory module to generate memory-embeddings that are then soft-prompt injected into an LLM. Mem0 (Chhikara et al. 2025) dynamically extracts, consolidates, and retrieves salient information from ongoing LLM conversations, with an additional graph-based variant. (Lee et al. 2024) use prompt engineering to dynamically compress and store memory content with optional referencing to the original textual context, while A-MEM (Xu et al. 2025) uses a more agentic approach by designing an agentic memory system that creates interconnected knowledge networks through dynamic indexing and linking. (Maharana et al. 2024) introduced a comprehensive pipeline to generate synthetic data over many conversation turns using different personas, showing persisting challenges in LLMs accurately processing lengthy dialogues. 


\subsection{LLMs in Mental Health}
As rapid advancements in LLMs continue, recent works have investigated their integration into a variety of domains, with one of them being mental health therapy. (Liu et al. 2023) introduced ChatCounselor, a LLaMA-7B model fine-tuned on psychology-based instruction pairs with performance nearing ChatGPT, showing the significant potential in supervised fine-tuning of LLMs for specific domains. Similarly, (Liu et al. 2025) introduced EmoScan, a multi-agentic system consisting of 2 fine-tuned Mistral-7B models, demonstrating performance improvements over base models and larger LLMs like GPT-4 in screening emotional disorders. BOLT (Chiu et al. 2024) was proposed as a comprehensive evaluation framework to assess LLMs’ current strengths and limitations in mental therapy, and found that LLMs tend to offer solutions over showing empathy when clients expressed their emotions. Similarly, (Zhang et al. 2025) introduced CBT-Bench to test LLMs in basic cognitive behavioral therapy knowledge acquisition, cognitive model understanding, and therapeutic response generation, with extensive experiments showing shortcomings of LLMs in complex real-world scenarios requiring deep patient analysis. 


\subsection{Test-Time Inference and Continual Learning}
Despite the massive amounts of data LLMs are trained on throughout the internet, there has been much ongoing research focused on further improving LLM performance without relying solely on said training data. The foundational chain-of-thought approach (Wei et al. 2023) improves mathematical and logical reasoning in LLMs by prompting them to output each reasoning step they use before stating a final answer, significantly reducing hallucinations and improving accuracy. More sophisticated approaches have been introduced as well, such as Titans (Behrouz et al. 2024) which introduces a neural long-term memory module that memorizes long historical context and assists the process of attending to current context while also looking at long past information. (Tandon et al. 2025) proposed an end-to-end pipeline approach for test-time training over long contexts by formulating it as a continual learning problem rather than an architecture design problem, while (Hu et al. 2025) instead dynamically adapts LLMs to specific domains using only unlabelled testing data for test-time training. Meanwhile (Snell et al. 2024) sought to determine how well LLMs can use additional but fixed test-time compute budget to improve performance, finding that the effectiveness of scaling up test-time compute critically depends on the difficulty of the prompt. Recently, (Behrouz et al. 2025) introduced nested learning as a new paradigm for machine learning in order to mitigate catastrophic forgetting in neural networks, by designing a multi-model system in which different models update at differing frequencies, analogous to the brain storing short-term and long-term memories. 


% \section{Methodology: Dynamic Test-Time Memory} -- Armaan
% \subsection{The Generator-Extractor Architecture}
% Explain the dual-component system:
% \begin{itemize}
%     \item \textbf{The Generator (The Summarizer):} Responsible for real-time synopsis creation. Detail the logic for creating synopses. Does it summarize per turn, per session, or per clinical breakthrough?
%     \item \textbf{The Extractor (The Retriever):} Handles context-aware retrieval of historical synopses. How does it identify relevant synopses based on the current user prompt?
% \end{itemize}

% \subsection{Inference-Time Continual Learning}
% Detail the recursive update loop that allows the memory state to evolve without weight updates. Describe the recursive nature and how the memory state evolves without retraining weights.

% \subsection{Prompt Engineering \& CBT Integration}
% How the system is ``steered" to maintain the therapist persona.

\section{Evaluation}

\subsection{Experimental Setup}

We evaluate on three synthetic patient transcripts (SYN\_0001, SYN\_0002, SYN\_0004), each generated via a Gemma-based patient--therapist simulation. All transcripts contain 50--78 total turns (25--39 counselor turns evaluated). We compare four conditions across two axes: (1)~who generates the counselor response (human vs.\ LLM), and (2)~what context the system receives (sliding window vs.\ accumulated memories).

A critical design choice: memory-based conditions receive \emph{only} retrieved memories with no sliding window, while context-based conditions receive \emph{only} the sliding window with no memories. This isolation tests whether each memory system can serve as the \textbf{sole} source of patient context, rather than merely supplementing a conversation window.

Table~\ref{tab:config} summarizes the input configuration for each condition.

\begin{table}[h]
\centering
\caption{Experimental conditions: what each counselor and judge receives as input. Memory-based conditions receive \emph{no} sliding window; context-based conditions receive \emph{no} memories.}
\label{tab:config}
\small
\begin{tabular}{lcccc}
\toprule
\textbf{Condition} & \textbf{Counselor} & \textbf{Counselor Input} & \textbf{Judge Input} \\
\midrule
Human (Context) & Human & N/A & Sliding window (10 turns) \\
Human (Memory) & Human & N/A & Mem0 memories only \\
LLM (Context Window) & LLM & Sliding window (10 turns) & Sliding window (10 turns) \\
LLM (Mem0) & LLM & Mem0 memories only & Mem0 memories only \\
DTM (Ours) & LLM & Curated cheatsheet only & \textit{TBD} \\
\bottomrule
\end{tabular}
\end{table}

\begin{itemize}
    \item \textbf{Human Counselor:} The original therapist responses from the synthetic transcript, serving as the \emph{upper bound} for therapeutic quality. We report two sub-conditions based on what the LLM judge receives during evaluation:
    \begin{itemize}
        \item \textbf{Human (Context):} The judge receives only the last 10 conversation turns (sliding window). No memory retrieval is performed.
        \item \textbf{Human (Memory):} The judge receives only Mem0's accumulated memories for the patient. No sliding window is provided.
    \end{itemize}
    Comparing these two sub-conditions reveals whether memory-aware evaluation introduces scoring bias on \emph{identical} counselor responses.
    \item \textbf{LLM (Context Window):} An LLM counselor that generates responses using the patient query and the last 10 conversation turns. The judge also receives only the sliding window. No memory system is used for either the counselor or the judge. This is the \emph{sliding window baseline}.
    \item \textbf{LLM (Mem0):} An LLM counselor that generates responses using \emph{only} Mem0's retrieved memories (no sliding window). The judge also receives only memories. This tests whether \emph{static memory accumulation} can replace conversational context.
    \item \textbf{DTM (Ours):} Our Dynamic Test-time Memory system using the DC-RS architecture. The counselor receives curated cheatsheet strategies instead of raw memories. \textit{(Results pending.)}
\end{itemize}

All conditions use the same judge model (\texttt{gpt-oss:20b}) and identical evaluation rubrics to ensure comparability.

\subsection{Evaluation Metrics}

We adopt two LLM-as-Judge metrics, each scored on a 1--10 scale per counselor turn:

\begin{itemize}
    \item \textbf{CBT Adherence Score (Part B -- Methodological Drift):} Measures how well the counselor follows Cognitive Behavioral Therapy principles. High scores (9--10) indicate excellent use of Socratic questioning, collaborative exploration, and evidence examination. Low scores (1--4) indicate directive advice-giving, ``should'' statements, or abandonment of CBT techniques.
    \item \textbf{Persona Consistency Score (Part C -- Boundary Dissolution):} Measures whether the counselor maintains a professional therapeutic persona. High scores indicate appropriate boundaries, professional language, and emotional distance. Low scores indicate mirroring of client slang, peer-like tone, or self-disclosure.
    \item \textbf{Alignment:} A composite measure of overall therapeutic quality, computed as the harmonic mean of CBT Adherence and Persona Consistency.
\end{itemize}

\subsection{Results}

\begin{table}[h]
\centering
\caption{Aggregated results across all three patients (SYN\_0001, SYN\_0002, SYN\_0004). Mean $\pm$ std.\ reported. $n$: total counselor turns evaluated. Best per-column in \textbf{bold}.}
\label{tab:aggregate}
\small
\begin{tabular}{lccccc}
\toprule
\textbf{Condition} & $n$ & \textbf{CBT Adh.} & \textbf{Persona Con.} & \textbf{Alignment} \\
\midrule
Human (Context) & 92 & 7.72 $\pm$ 1.39 & \textbf{9.02 $\pm$ 0.61} & 8.32 \\
Human (Memory) & 92 & 7.29 $\pm$ 1.55 & 8.86 $\pm$ 0.46 & 8.00 \\
LLM (Context Window) & 93 & \textbf{8.57 $\pm$ 1.98} & 8.65 $\pm$ 1.70 & \textbf{8.61} \\
LLM (Mem0) & 93 & 7.67 $\pm$ 2.91 & 7.89 $\pm$ 2.66 & 7.78 \\
DTM (Ours) & --- & \multicolumn{3}{c}{\textit{Results pending}} \\
\bottomrule
\end{tabular}
\end{table}

\begin{table}[h]
\centering
\caption{Per-patient breakdown. Mean $\pm$ std.\ reported for each metric.}
\label{tab:perpatient}
\small
\begin{tabular}{llccccc}
\toprule
\textbf{Patient} & \textbf{Condition} & $n$ & \textbf{CBT Adh.} & \textbf{Persona Con.} & \textbf{Align.} \\
\midrule
\multirow{4}{*}{SYN\_0001}
 & Human (Context) & 28 & 7.46 $\pm$ 1.91 & 9.07 $\pm$ 0.90 & 8.19 \\
 & Human (Memory) & 28 & 7.00 $\pm$ 2.04 & 8.89 $\pm$ 0.50 & 7.83 \\
 & LLM (Context Window) & 29 & \textbf{9.10 $\pm$ 0.56} & \textbf{9.03 $\pm$ 0.33} & \textbf{9.07} \\
 & LLM (Mem0) & 29 & 7.97 $\pm$ 2.54 & 8.41 $\pm$ 2.08 & 8.18 \\
\midrule
\multirow{4}{*}{SYN\_0002}
 & Human (Context) & 39 & 7.92 $\pm$ 1.11 & \textbf{9.00 $\pm$ 0.46} & 8.43 \\
 & Human (Memory) & 39 & 7.56 $\pm$ 1.21 & 8.77 $\pm$ 0.48 & 8.12 \\
 & LLM (Context Window) & 39 & \textbf{8.26 $\pm$ 2.37} & 8.51 $\pm$ 1.90 & \textbf{8.38} \\
 & LLM (Mem0) & 39 & 7.49 $\pm$ 3.16 & 7.64 $\pm$ 2.84 & 7.56 \\
\midrule
\multirow{4}{*}{SYN\_0004}
 & Human (Context) & 25 & 7.68 $\pm$ 1.03 & \textbf{9.00 $\pm$ 0.41} & 8.29 \\
 & Human (Memory) & 25 & 7.20 $\pm$ 1.38 & 8.96 $\pm$ 0.35 & 7.98 \\
 & LLM (Context Window) & 25 & \textbf{8.44 $\pm$ 2.29} & 8.40 $\pm$ 2.24 & \textbf{8.42} \\
 & LLM (Mem0) & 25 & 7.60 $\pm$ 3.00 & 7.68 $\pm$ 2.98 & 7.64 \\
\bottomrule
\end{tabular}
\end{table}

\subsection{Comparison Against Baselines}

Aggregated results across all three patients are reported in Table~\ref{tab:aggregate}.

\paragraph{LLM (Context Window) achieves the highest overall alignment (8.61).} With a CBT adherence of 8.57 and persona consistency of 8.65, it outperforms all other conditions. This is expected: when the LLM sees only the last 10 turns, it remains tightly focused on the immediate therapeutic exchange without distraction from noisy historical context. However, this advantage comes at the cost of \emph{long-range continuity}---the model cannot reference patient history beyond its window.

\paragraph{Human counselor responses set the persona ceiling.} Under context-only evaluation, the original therapist responses score 9.02 on persona consistency---the highest of any condition---reflecting the natural professional tone of a trained counselor. CBT adherence is lower (7.72), likely because the Gemma-simulated therapist occasionally uses directive language or omits Socratic questioning in closing turns.

\paragraph{Memory-aware judging slightly deflates scores on identical responses.} Comparing the two Human sub-conditions reveals a consistent pattern: when the judge receives accumulated memories instead of a sliding window, scores decrease (CBT: 7.72 $\to$ 7.29; Persona: 9.02 $\to$ 8.86; Alignment: 8.32 $\to$ 8.00). This suggests that raw Mem0 memories introduce evaluative noise---the judge may penalize responses that appear inconsistent with noisy or redundant memory entries, even when the responses themselves are unchanged. This finding underscores the importance of \emph{memory curation}: uncurated memories can distort not only generation but also downstream evaluation.

\paragraph{LLM (Mem0) scores lowest across all metrics.} With CBT adherence of 7.67, persona consistency of 7.89, and alignment of 7.78, the Mem0-augmented LLM performs worse than both the context window baseline and the human counselor. Critically, it also exhibits the highest variance (standard deviations up to 3.16 for CBT and 2.98 for persona), indicating \emph{unpredictable performance}. This supports our core hypothesis: na\"ive memory accumulation without curation leads to distortion buildup that harms therapeutic alignment over extended conversations.

\subsection{Per-Patient Breakdown}

The per-patient results in Table~\ref{tab:perpatient} reveal several important patterns.

\paragraph{Longer transcripts amplify Mem0 degradation.} SYN\_0002, the longest transcript (39 evaluated turns), shows the most pronounced degradation under Mem0: CBT drops to 7.49 $\pm$ 3.16, and persona falls to 7.64 $\pm$ 2.84. In contrast, both Human (Context) and LLM (Context Window) remain stable on the same transcript. This is consistent with the \emph{memory distortion accumulation} problem: as more memories are stored without curation, retrieval quality degrades over longer conversations.

\paragraph{Context Window excels on shorter transcripts.} On SYN\_0001, the Context Window baseline achieves a CBT score of 9.10 $\pm$ 0.56 and alignment of 9.07---the tightest distribution and highest alignment of any condition on any patient. Its advantage narrows on longer transcripts (alignment drops to 8.38 on SYN\_0002), hinting that even the sliding window approach faces challenges as conversations extend. This motivates the need for a memory system that can maintain quality over longer horizons.

\paragraph{Human persona consistency is remarkably stable.} Across all three patients under context-only judging, the Human Counselor maintains a persona score between 9.00--9.07 with standard deviations below 0.90. This near-ceiling performance with minimal variance sets the target that any memory-augmented LLM system must approach.

\paragraph{Memory-aware judging consistently penalizes the same responses.} The gap between Human (Context) and Human (Memory) is consistent across patients: alignment drops by 0.36 on SYN\_0001, 0.31 on SYN\_0002, and 0.31 on SYN\_0004. This systematic effect suggests that the memory-induced evaluation bias is not patient-specific but a property of how raw memories interact with the judge's scoring process.

\subsection{Discussion}

These baseline results establish three key findings:

\begin{enumerate}
    \item \textbf{CBT adherence vs.\ long-range memory:} The Context Window baseline excels at CBT adherence and alignment precisely because it avoids memory noise, but it cannot maintain continuity across sessions. A successful memory system must match its per-turn scores while adding cross-session awareness. The degradation of Context Window performance on longer transcripts (alignment: 9.07 on SYN\_0001 $\to$ 8.38 on SYN\_0002) further motivates this need.

    \item \textbf{Raw memory hurts both generation and evaluation:} Mem0 provides the counselor with patient history, yet this comes at the cost of both methodological drift and persona dissolution. Moreover, when the \emph{same} human responses are evaluated with raw memories as context, scores drop compared to sliding-window evaluation. The uncurated nature of Mem0's stored memories introduces noise that degrades quality on both sides of the pipeline---exactly the failure mode DTM is designed to address through strategy curation and clinical filtering.

    \item \textbf{Variance as a reliability signal:} Beyond mean scores, variance across turns is a critical indicator of system reliability. LLM (Mem0) exhibits standard deviations 1.5--2$\times$ larger than other conditions, meaning a therapist using this system would experience highly unpredictable turn-by-turn quality---unacceptable in a clinical setting where consistency is paramount.
\end{enumerate}

DTM (Ours) aims to resolve these tensions by replacing raw memory accumulation with curated therapeutic strategies that preserve CBT alignment while enabling long-range patient context. Results are forthcoming upon completion of the \texttt{therapy\_dc\_rs} evaluation pipeline.

\section{Experimental Setup}

\subsection{Dataset} -- Daniel
Description of curated multi-turn, multi-session therapy transcripts.

\subsection{Baselines}
Baselines: 1. LLM + Standard History (Context Window). 2. LLM + Mem0 (External Memory).

\subsection{Evaluation Metrics}
\begin{itemize}
    \item \textbf{CBT Score:} Specialized rubric for therapeutic adherence.
    \item \textbf{Persona Consistency:} Stability of the counselor's identity.
    \item \textbf{Gemini-as-a-Judge:} Methodology for using frontier models as evaluators.
\end{itemize}

\section{Results and Discussion}

\subsection{Performance Comparison}
Showing the ``Degradation Curve" on how standard LLMs fail as session count increases vs. the stability of our DTM approach.
\subsection{Comparison Against Baselines}

\begin{table}[h]
\centering
\caption{Comparison of DTM against standard history and Mem0 baselines. Scores are averaged across all three synthetic patients (SYN\_0001, SYN\_0002, SYN\_0004) on a 1--10 scale.}
\begin{tabular}{@{}lccc@{}}
\toprule
\textbf{Method} & \textbf{CBT Score} & \textbf{Consistency} & \textbf{Alignment} \\ \midrule
Human Counselor          & 7.69 & 9.02 & -- \\
LLM (Context Window)     & 8.60 & 8.65 & -- \\
LLM (Mem0)               & 7.68 & 7.91 & -- \\
\textbf{DTM (Ours)}      & --   & --   & -- \\ \bottomrule
\end{tabular}
\end{table}

\begin{table}[h]
\centering
\caption{Per-patient breakdown of CBT Adherence and Persona Consistency scores (mean $\pm$ std).}
\resizebox{\textwidth}{!}{%
\begin{tabular}{@{}lcccccc@{}}
\toprule
& \multicolumn{2}{c}{\textbf{SYN\_0001}} & \multicolumn{2}{c}{\textbf{SYN\_0002}} & \multicolumn{2}{c}{\textbf{SYN\_0004}} \\
\cmidrule(lr){2-3} \cmidrule(lr){4-5} \cmidrule(lr){6-7}
\textbf{Method} & \textbf{CBT} & \textbf{Consist.} & \textbf{CBT} & \textbf{Consist.} & \textbf{CBT} & \textbf{Consist.} \\ \midrule
Human Counselor      & 7.46 {\scriptsize$\pm$1.91} & 9.07 {\scriptsize$\pm$0.90} & 7.92 {\scriptsize$\pm$1.11} & 9.00 {\scriptsize$\pm$0.46} & 7.68 {\scriptsize$\pm$1.03} & 9.00 {\scriptsize$\pm$0.41} \\
LLM (Context Window) & \textbf{9.10} {\scriptsize$\pm$0.56} & \textbf{9.03} {\scriptsize$\pm$0.33} & \textbf{8.26} {\scriptsize$\pm$2.37} & 8.51 {\scriptsize$\pm$1.90} & \textbf{8.44} {\scriptsize$\pm$2.29} & 8.40 {\scriptsize$\pm$2.24} \\
LLM (Mem0)           & 7.97 {\scriptsize$\pm$2.54} & 8.41 {\scriptsize$\pm$2.08} & 7.49 {\scriptsize$\pm$3.16} & 7.64 {\scriptsize$\pm$2.84} & 7.60 {\scriptsize$\pm$3.00} & 7.68 {\scriptsize$\pm$2.98} \\
\textbf{DTM (Ours)}  & --   & --   & --   & --   & --   & --   \\ \bottomrule
\end{tabular}%
}
\end{table}
\subsection{Turn-wise evaluation breakdown}


\subsection{Qualitative Analysis}
Case studies highlighting how the Extractor maintains long-term clinical threads. Especially show cases of ``Memory Hits" where the Extractor pulled a crucial detail from Session 1 to solve a problem in Session 5 for example.

\subsection{Inference Latency}
A brief look at the trade-off between memory ``reflection" time and response quality.

\section{Conclusion}
Summary of findings and implications for the deployment of AI agents in "the wild" clinical settings.

\subsubsection*{Author Contributions}
If you'd like to, you may include  a section for author contributions as is done
in many journals. This is optional and at the discretion of the authors.

\subsubsection*{Acknowledgments}
Use unnumbered third level headings for the acknowledgments. All
acknowledgments, including those to funding agencies, go at the end of the paper.


\bibliography{iclr2026_conference}
\bibliographystyle{iclr2026_conference}

\appendix
\section{Appendix}
You may include other additional sections here.


\end{document}
